% 太阳系（综述）
% license CCBYSA3
% type Wiki

本文根据 CC-BY-SA 协议转载翻译自维基百科\href{https://en.wikipedia.org/wiki/Solar_System}{相关文章}。

\begin{figure}[ht]
\centering
\includegraphics[width=8cm]{./figures/5a727e2db62d7a0e.png}
\caption{太阳、行星、卫星与矮行星\(^\text{[a]}\)（真实色彩，比例为尺寸等比，但距离不按比例）} \label{fig_Solar_1}
\end{figure}
太阳系（The Solar System）由太阳及其所环绕的天体组成。该名称来源于“Sōl”，即太阳的拉丁名称。大约在 46 亿年前，一个分子云中密度较高的区域发生坍缩，形成了太阳以及一个原行星盘（protoplanetary disc），环绕天体便由此汇聚而成。太阳核心内部的氢融合为氦，释放出能量，这些能量主要通过其外层光球（photosphere）向外辐射，从而在整个系统中形成一个由内而外逐渐降低的温度梯度。太阳系中超过 99.86\% 的质量集中在太阳内部。

在绕太阳运行的天体中，质量最大的为八大行星。按与太阳距离递增的顺序，最靠近太阳的四颗类地行星（terrestrial planets）分别是—— 水星（Mercury）、金星（Venus）、地球（Earth）和火星（Mars）。它们构成了太阳系的内行星系统（inner Solar System）。其中，地球与火星是太阳系中仅有的两颗位于太阳宜居带（habitable zone）内的行星——在此范围内，行星表面可存在液态水。

在距离太阳约 5 个天文单位（AU）之外、越过“霜线”（frost line）的位置，分布着四颗更为庞大的行星：两颗气态巨行星（gas giants）—— 木星（Jupiter）与土星（Saturn），以及两颗冰巨行星（ice giants）—— 天王星（Uranus）与 海王星（Neptune）。它们共同构成了太阳系的外行星系统（outer Solar System）。其中，木星与土星占据了太阳系除恒星外几乎 90\% 的全部质量。

太阳系中还存在着数量庞大但质量较小的天体。天文学家普遍认为，太阳系至少拥有九颗矮行星（dwarf planets），分别是：谷神星（Ceres）、奥库斯（Orcus）、冥王星（Pluto）、哈梅亚（Haumea）、夸奥尔（Quaoar）、鸟神星（Makemake）、共工（Gonggong）、阋神星（Eris）与赛德娜（Sedna）。在这些天体中，有六颗行星、七颗矮行星及若干其他天体拥有自然卫星（natural satellites），即通常所称的“月亮（moons）”。这些卫星的大小范围极广：从如地球之月那样接近矮行星尺度的巨型卫星，到仅有极小质量的“月粒（moonlets）”。此外，太阳系中还存在大量小型天体（small Solar System bodies），包括小行星（asteroids）、彗星（comets）、半人马天体（centaurs）、流星体（meteoroids）以及星际尘埃云（interplanetary dust clouds）。其中一部分天体分布在小行星带（asteroid belt）（位于火星与木星轨道之间）以及柯伊伯带（Kuiper belt）（位于海王星轨道之外）。

在这些天体之间，存在着充满尘埃与粒子的行星际介质（interplanetary medium）。太阳不断释放的带电粒子形成太阳风（solar wind），它向外扩散，构成一个被称为日球层（heliosphere）的巨大空间。在距离太阳约 70–90 个天文单位（AU）处，太阳风受到星际介质（interstellar medium）的阻挡，形成边界——日球顶（heliopause）。这一区域标志着太阳系与星际空间的分界。更远处、约 2,000 个天文单位（AU）以外，是太阳系最外层的假想区域——奥尔特云（Oort cloud）。它被认为是长周期彗星（long-period comets）的来源，其范围可能延伸至太阳系的引力界限（Hill sphere）边缘，约178,000–227,000 AU（2.81–3.59 光年），在此处太阳系的引力势与银河系的引力势达到平衡。目前，太阳系正穿越一个被称为本地星际云（Local Cloud）的星际介质区域。距离太阳系最近的恒星是比邻星（Proxima Centauri），距离约269,000 AU（4.25 光年）。太阳系与比邻星都位于本地泡（Local Bubble）内，这是银河系中一个直径约 1,000 光年（ly）的相对稀薄的区域。
\subsection{定义}
太阳系包括太阳以及所有受其引力约束并围绕其运行的天体。\(^\text{[15][16][17]}\)

国际天文学联合会（IAU）将太阳系描述为：由太阳引力所束缚的一切天体，包括太阳本身、八大行星，以及环绕太阳运行的其他天体。\(^\text{[11]}\)美国国家航空航天局（NASA）则将太阳系定义为一个行星系统（planetary system），其中包括太阳及所有环绕它运行的天体。\(^\text{[12]}\)

当“solar system”一词不作为专有名词、且首字母不大写时，它可以指太阳系本身，也可泛指任何类似太阳系的行星系统。\(^\text{[15]}\)
\subsection{形成与演化}
\subsubsection{过去}
\begin{figure}[ht]
\centering
\includegraphics[width=8cm]{./figures/0dc12bd0e5f767c3.png}
\caption{} \label{fig_Solar_2}
\end{figure}
太阳系形成于至少约46.68 亿年前，源自一片巨大分子云中某个区域的引力坍缩。\(^\text{[b]}\)这片最初的云可能横跨数光年，并极有可能同时孕育了多颗恒星。\(^\text{[19]}\)如同其他分子云一样，它的主要成分是氢气，其次是氦气，并含有少量由前几代恒星通过核聚变产生的较重元素。\(^\text{[20]}\)

当原太阳星云\(^\text{[20]}\)开始坍缩时，角动量守恒使其旋转速度不断加快。中心区域由于聚集了大部分质量，变得比周围更炽热、更致密。\(^\text{[16]}\)随着收缩的星云自转加速，它逐渐扁平化为一个直径约 200 天文单位（AU）的原行星盘\(^\text{[19][21]}\)，其中心形成了一个炽热而致密的原恒星。\(^\text{[22][23]}\)行星由盘中的物质通过吸积（accretion）过程形成，\(^\text{[24]}\)其中尘埃和气体相互引力吸引，逐渐聚合成更大的天体。在早期太阳系中，可能存在数百个原行星（protoplanets），但它们大多互相合并、被摧毁或被抛射出系统，最终仅留下行星、矮行星以及少量残余小天体。\(^\text{[25][26]}\)

在靠近太阳的内太阳系，即霜线（frost line）以内、甚至烟尘线（soot line）以内，\(^\text{[27]}\)由于高温使除金属与硅酸盐外的物质难以以固体形式存在，因此这一带形成的行星主要由岩石构成，即：水星、金星、地球与火星。由于这些耐高温的物质仅占太阳星云极小的比例，类地行星因此无法增长得太大。\(^\text{[25]}\)

四颗巨行星——木星、土星、天王星与海王星——在霜线之外形成，该区域位于火星与木星轨道之间，温度足够低，使挥发性冰类化合物可以保持固态。这些冰的丰度远高于类地行星形成所用的金属与硅酸盐，这使得巨行星能长得足够大，从而捕获大量的氢与氦 —— 这两种元素是宇宙中最轻、最丰富的。\(^\text{[25]}\)未能成为行星的残余碎片则聚集在小行星带、柯伊伯带与奥尔特云等区域。\(^\text{[25]}\)

约五千万年后，原恒星中心的氢气压与密度足够高，开始发生热核聚变反应。\(^\text{[28]}\)随着核心中氦的积累，太阳逐渐变得更亮；\(^\text{[29]}\)在主序星生命早期，其亮度仅为如今的约 70\%。\(^\text{[30]}\)温度、反应速率、压力与密度不断上升，直至达到静水平衡（hydrostatic equilibrium）：即热压与引力达到平衡。此时，太阳成为一颗主序星。\(^\text{[31]}\)太阳风从太阳表面喷射，形成了日球层（heliosphere），并将原行星盘中剩余的气体与尘埃吹散至星际空间。\(^\text{[29]}\)

在原行星盘消散后，尼斯模型（Nice model）提出：气体巨行星与微行星之间的引力相互作用导致行星迁移至不同轨道。这引发了整个系统的动力学不稳定，使微行星被散射，最终将气体巨行星带入如今的位置。在这一时期，大调车假说（Grand Tack Hypothesis）认为：木星的一次最终向内迁移扰乱了小行星带的大部分物质，从而引发了晚期重轰炸（Late Heavy Bombardment）。\(^\text{[32][33]}\)
\subsubsection{现在与未来}
太阳系目前处于一个相对稳定、缓慢演化的状态，所有天体都沿着受引力束缚的轨道绕太阳运行。\(^\text{[34]}\)尽管太阳系在数十亿年来一直较为稳定，但从技术角度上看，它是一个混沌系统，未来可能受到扰动。在未来数十亿年内，存在极小概率会有其他恒星穿越太阳系。这可能导致系统的轻微失稳，在数百万年后造成行星被逐出、相互碰撞或坠入太阳，但最有可能的结果仍是：太阳系将大体保持如今的样貌。\(^\text{[35]}\)
\begin{figure}[ht]
\centering
\includegraphics[width=8cm]{./figures/080701ee703e48af.png}
\caption{当前太阳与其在红巨星阶段达到峰值时的体积对比} \label{fig_Solar_3}
\end{figure}
太阳的主序阶段（main-sequence phase）——从开始到结束——将持续约 100 亿年，而此后所有阶段（从主序末期到白矮星前的各阶段）加起来的寿命仅约 20 亿年。\(^\text{[36]}\)太阳系将在此期间大体保持如今的状态，直到太阳核心中的氢完全转化为氦为止—— 这将发生在大约 50 亿年后。届时，太阳的主序生命将走到尽头。  

那时，太阳的核心将因引力进一步收缩，氢聚变将仅在包围惰性氦核的壳层（shell）中进行，其总能量输出将超过当前水平。太阳的外层将剧烈膨胀，直径将达到目前的约260 倍，并演化为一颗红巨星（Red Giant）。由于表面积大幅增加，其表面温度将下降，在最冷时约为2,600 K（4,220 °F），远低于主序阶段的温度。\(^\text{[36]}\)

膨胀后的太阳预计将汽化水星与金星，并使地球与火星变得不可居住（甚至可能完全摧毁地球）。\(^\text{[37][38]}\)最终，太阳核心温度将升高到足以点燃氦聚变反应；太阳将在核心燃烧氦，但这一阶段仅持续其氢燃烧阶段的极小一部分时间。由于太阳的质量不足以启动更重元素的聚变，核心中的核反应将逐渐衰竭并停止。太阳的外层将在这一过程中被抛射到太空中，留下一个致密的白矮星（White Dwarf）—— 其质量约为太阳原始质量的一半，但体积仅有地球大小。\(^\text{[36]}\)被抛射的外层气体可能形成一个行星状星云（Planetary Nebula），将太阳形成时的部分物质—— 此时已被碳等重元素所“富集”—— 重新送回星际介质（Interstellar Medium）。\(^\text{[39][40]}\)
\subsection{一般特征}
\begin{figure}[ht]
\centering
\includegraphics[width=8cm]{./figures/7b113d47a833549c.png}
\caption{这是一张经过色彩增强的从月球拍摄的太阳系多种成分的照片。左下角的三个光点自左向右依次为行星——土星、火星与水星；画面中央，太阳的日冕（corona）从月球漆黑的边缘上方升起，而月球右侧被地球反照光（earthshine）所照亮。} \label{fig_Solar_4}
\end{figure}
天文学家有时会将太阳系的结构划分为若干独立的区域。内太阳系（Inner Solar System）包括水星、金星、地球、火星以及位于小行星带中的天体；外太阳系（Outer Solar System）包括木星、土星、天王星、海王星以及位于柯伊伯带中的天体。\(^\text{[41]}\)自从发现柯伊伯带以来，太阳系的最外层部分被视为一个独立的区域，由位于海王星之外的天体组成。\(^\text{[42]}\)
\subsubsection{组成}
太阳系的主要成分是太阳，一颗G型主序星（G-type main-sequence star），其质量占整个系统已知总质量的99.86\%，在引力上完全主导着整个系统。\(^\text{[43]}\)围绕太阳运行的四个最大天体——四大行星——占据剩余质量的99\%，其中仅木星和土星就合计超过90\%。太阳系的其他所有天体（包括四颗类地行星、矮行星、卫星、小行星与彗星等）总共的质量还不到太阳系总质量的0.002\%。\(^\text{[h]}\)

太阳的成分大约由 98\%的氢与氦组成，\(^\text{[47]}\)木星与土星的成分也大体相同。\(^\text{[48][49]}\)太阳系中存在一个由早期太阳的热量与光压形成的成分梯度：距离太阳较近、受热与光压影响更强的天体，主要由熔点较高的元素构成；而距离太阳较远的天体，则主要由熔点较低的物质组成。\(^\text{[50]}\)在太阳系中，那条易挥发物质能够凝聚的界限称为霜线（frost line），其位置大约在距离太阳约地球距离的五倍处。\(^\text{[5]}\)
\subsubsection{轨道}
行星与其他围绕太阳运行的大型天体的轨道大多位于太阳系的不变平面（invariable plane）附近，地球的轨道（称为黄道面，ecliptic）也接近这一平面；其中木星的轨道与该平面的夹角最小，仅为 0.3219°。\(^\text{[51]}\)而较小的冰质天体（例如彗星）则常以更大的倾角绕此平面运行。\(^\text{[52][53]}\)太阳系中的大多数行星都拥有自己的“次级系统”，即被称为卫星（moons）的天然卫星所环绕。所有最大型的天然卫星都呈同步自转（synchronous rotation），始终以一面对着其母行星。四大行星（木星、土星、天王星、海王星）皆拥有行星环（planetary rings）—— 由无数微小颗粒组成的薄环盘，这些颗粒共同绕行行星。\(^\text{[54]}\)

由于太阳系形成的原因，行星与绝大多数其他天体的公转方向与太阳自转方向一致，即从地球北极上方观察为逆时针方向（counter-clockwise）。\(^\text{[55]}\)不过也有例外，例如哈雷彗星（Halley’s Comet）。\(^\text{[56]}\)多数较大的卫星以顺行方向（prograde direction）绕行行星，即与母行星的自转方向一致；而海王星的卫星海卫一（Triton）是已知最大的一颗逆行（retrograde）运行的卫星。\(^\text{[57]}\)大多数大型天体的自转方向相对于其公转方向也是顺行的，但金星（Venus）的自转方向为逆行。\(^\text{[58]}\)

在良好的近似下，开普勒行星运动定律（Kepler’s laws of planetary motion）可以描述天体绕太阳的轨道运动。\(^\text{[59]: 433–437}\) 这些定律指出：每个天体沿着一个椭圆轨道运行，而太阳位于该椭圆的一个焦点上，这导致天体在公转过程中与太阳的距离不断变化。天体距太阳最近的位置称为近日点（perihelion），最远的位置称为远日点（aphelion）。\(^\text{[60]: 9–6 }\)除水星外，行星的轨道几乎为圆形，而许多彗星、小行星和柯伊伯带天体则具有高度椭圆的轨道。开普勒定律只考虑太阳对轨道天体的引力作用，未考虑天体之间的相互引力。在人类的时间尺度上，这些扰动可通过数值模型加以修正，\(^\text{[60]: 9–6}\) 但在数十亿年的时间尺度上，整个行星系统可能会发生混沌变化。\(^\text{[61]}\)

太阳系的角动量（angular momentum）衡量了系统内所有运动成分的总轨道与自转动量。\(^\text{[62]}\)虽然太阳的质量占据绝对主导，但其仅贡献了大约2\% 的角动量。\(^\text{[63][64]}\)行星——主要是木星——由于其质量、轨道与距离太阳的组合，占据了系统绝大部分的角动量，而彗星可能也提供了显著的次要贡献。\(^\text{[63]}\)
\subsubsection{距离与尺度}
\begin{figure}[ht]
\centering
\includegraphics[width=8cm]{./figures/5dbc152db58202da.png}
\caption{太阳系中各行星相对轨道距离的可视化示意图——以压缩矩形形式呈现。} \label{fig_Solar_5}
\end{figure}
太阳的半径为0.0047 AU（70万千米；40万英里）。\(^\text{[65]}\)因此，太阳的体积仅占一个半径等于地球轨道半径的球体体积的0.00001\%（即 1/10^7），而地球的体积大约仅为太阳的百万分之一（10^{-6}）。最大行星木星距离太阳5.2 AU，半径为71,000 千米（0.00047 AU；44,000 英里）；最远的行星海王星则距离太阳30 AU。\(^\text{[49][66]}\)

除少数例外外，行星或行星带距离太阳越远，其轨道间距也越大。例如：金星比水星远约 0.33 AU；而土星比木星远4.3 AU；海王星则比天王星远10.5 AU。天文学家曾尝试寻找这些轨道间距的规律，如提丢斯–波得定律（Titius–Bode law）\(^\text{[67]}\)以及开普勒（Johannes Kepler）基于柏拉图立体（Platonic solids）的模型\(^\text{[68]}\)，但随着新的天体被发现，这些假说已被证伪。\(^\text{[69]}\)

太阳系比例模型（Solar System Scale Models）一些太阳系模型尝试以人类直观方式表现其中的相对尺度。 有的模型较小（甚至可为机械式，称为行星仪 orrery），也有的规模宏大，跨越整个城市或地区。\(^\text{[70]}\)目前已知最大的比例模型是瑞典的“瑞典太阳系模型（Sweden Solar System）”，它以斯德哥尔摩的阿维奇竞技场（Avicii Arena）—— 一座直径 110 米（361 英尺）的穹顶建筑——代表太阳。按照比例：木星是一颗直径 7.5 米（25 英尺）的球体，位于斯德哥尔摩阿兰达机场（40 千米 / 25 英里 之外）；而最远的天体 —— 赛德娜（Sedna）—— 则是一颗直径仅 10 厘米（4 英寸）的球体，位于吕勒奥（Luleå），距“太阳” 912 千米（567 英里）。\(^\text{[71][72]}\)在这一比例下，最近恒星比邻星（Proxima Centauri）的距离大约是地月距离的八倍。

若将太阳至海王星的距离缩放为100 米（330 英尺），则太阳的直径约为 3 厘米（1.2 英寸）—— 约为高尔夫球直径的三分之二；四大行星的直径均不超过 3 毫米（0.12 英寸）；而地球及其他类地行星的直径在该比例下将小于跳蚤（约 0.3 毫米 / 0.012 英寸）。\(^\text{[73]}\)
\begin{figure}[ht]
\centering
\includegraphics[width=8cm]{./figures/57f34943e92d54c3.png}
\caption{行星之间距离的比较图，白色横条表示轨道变化范围。行星的大小未按比例绘制。} \label{fig_Solar_6}
\end{figure}
\subsubsection{宜居性}
\begin{figure}[ht]
\centering
\includegraphics[width=8cm]{./figures/cc03a791c33b51ee.png}
\caption{太阳系与 TRAPPIST-1 恒星系宜居带的比较图。TRAPPIST-1 是一颗超冷红矮星，已知其周围有七颗类地行星以稳定轨道环绕运行。} \label{fig_Solar_7}
\end{figure}
\begin{figure}[ht]
\centering
\includegraphics[width=8cm]{./figures/caa7fcdf1175a2ac.png}
\caption{不同恒星温度下宜居带的比较图，示例包括若干已知系外行星以及地球、火星和金星。从上到下依次为：F 型主序星、黄矮星（G 型主序星）、橙矮星（K 型主序星）、典型红矮星，以及超冷矮星。} \label{fig_Solar_8}
\end{figure}
太阳系的宜居带（habitability zone）通常被认为位于内太阳系靠近地球的区域，在那里由于太阳辐射的作用，大气层中的液态水可以稳定存在。\(^\text{[74]}\)

除了太阳能外，太阳系中能够支持生命存在的主要特征是日球层（heliosphere）以及各行星自身的磁场（对于那些拥有磁场的行星而言）。这些磁场在一定程度上屏蔽了来自星际空间的高能粒子——宇宙射线。星际介质中宇宙射线的密度以及太阳磁场的强度会在极长的时间尺度上变化，因此太阳系中宇宙射线的渗透水平也随之变化，但其变化幅度目前仍未知。\(^\text{[75]}\)

太阳系中的宜居性不仅仅取决于表面条件或太阳辐射环境，因为在若干太阳系天体中可能存在的地下海洋中也可能具备生命存在的潜力，\(^\text{[76]}\)甚至在某些行星（尤其是金星）的云层中，也可能存在可居住层。\(^\text{[77]}\)
\subsubsection{与太阳系外行星系统的比较}
对开普勒（Kepler）望远镜观测数据的分析表明，银河系中被观测到的行星系统大致可分为三类：“相似型（similar）”系统——行星大小相近、间距相似且轨道高度圆形；“有序型（ordered）”系统——行星的质量随着与恒星距离的增加而增大；“混合型（mixed）”系统——行星质量分布无明显规律。太阳系属于“有序型”系统，占已观测系统的约37\%。然而，“相似型”系统为多数，占59\%，而“混合型”系统仅占4\%。\(^\text{[78]}\)

与许多太阳系外行星系统相比，太阳系的一个显著特征是 —— 缺乏位于水星轨道以内的行星。\(^\text{[79][80]}\)太阳系中也没有“超级地球”（质量为地球的1至10倍的行星），\(^\text{[79]}\)尽管假设中的“第九行星”（Planet Nine）若存在，可能是一颗位于太阳系边缘的超级地球。\(^\text{[81]}\)

另一个不寻常之处在于太阳系仅包含较小的类地行星和巨大的气体行星，而在其他行星系统中，中等大小的行星更为常见——包括岩质型与气态型——因此在地球与海王星（半径约为地球的3.8倍）之间的“尺寸空档”在太阳系中非常明显。由于许多超级地球的轨道都比水星更靠近其母星，有一种假说认为所有行星系统最初都有许多靠近恒星的行星，而随后它们通过一系列碰撞事件合并成更少但更大的行星。然而在太阳系中，这些碰撞反而导致行星的破坏与被抛出系统。\(^\text{[79][82]}\)

太阳系行星的轨道几乎为圆形。与许多其他系统相比，它们的轨道偏心率要小得多。\(^\text{[79]}\)尽管有尝试将这种现象部分归因于径向速度法的观测偏差，以及长期多行星相互作用的动力学效应，但其确切原因仍未确定。\(^\text{[79][83]}\)
\subsubsection{太阳}
\begin{figure}[ht]
\centering
\includegraphics[width=8cm]{./figures/cc5da4e20b80c9b8.png}
\caption{真实白光下的太阳（The Sun in true white color）} \label{fig_Solar_9}
\end{figure}
太阳（The Sun）是太阳系的中心恒星，也是其最为巨大的组成部分。它拥有巨大的质量（约为地球质量的 332,900 倍），\(^\text{[84]}\)占太阳系总质量的 99.86\%，\(^\text{[85]}\)这一庞大的引力使其核心达到足以维持氢聚变为氦的高温高压。\(^\text{[86]}\)这一聚变过程释放出极其巨大的能量，主要以电磁辐射的形式向外辐射，其辐射峰值位于可见光波段。\(^\text{[87][88]}\)

由于太阳在核心中进行氢聚变，因此它是一颗主序星。更具体地说，它属于G2 型主序星（G2V star），其中类型编号指其有效温度。主序星温度越高，其光度越强，但寿命越短。太阳的温度介于最热恒星与最冷恒星之间。比太阳更明亮、更炽热的恒星相对罕见；而显著更暗、更冷的恒星——即所谓的红矮星——约占银河系中所有核聚变恒星的 75\%。\(^\text{[89]}\)

太阳属于第一族恒星（Population I star），形成于银河系的旋臂中。与银河系核球与晕中的老年第二族恒星相比，它含有更多比氢与氦更重的元素（在天文学中称为“金属”）。\(^\text{[90]}\)这些较重元素是在更早一代恒星的核心与爆炸中形成的，因此宇宙必须先经历第一代恒星的死亡，才能被这些重元素所富集。最古老的恒星几乎不含金属，而后期形成的恒星则金属丰度更高。人们认为这种较高的金属丰度对太阳形成行星系统至关重要，因为行星正是由这些“金属”物质的吸积所形成的。\(^\text{[91]}\)

由太阳磁层主导的空间区域称为日球层（heliosphere），它覆盖了太阳系的大部分区域。除了光之外，太阳还不断释放带电粒子流（等离子体），称为太阳风（solar wind）。太阳风以每小时 90 万公里（56 万英里）至 288 万公里（179 万英里）的速度向外扩散，\(^\text{[92]}\)充满了太阳系天体之间的真空空间，形成了一种稀薄、含尘的介质，称为行星际介质（interplanetary medium），其范围至少延伸至 100 天文单位。\(^\text{[93]}\)

太阳表面的活动，如太阳耀斑（solar flares）与日冕物质抛射（coronal mass ejections, CME），会扰动整个日球层，产生 空间天气（space weather）并引发 地磁风暴（geomagnetic storms）。\(^\text{[94]}\)日冕物质抛射及类似事件会将磁场与大量物质从太阳表面喷射出来。这些磁场与物质与地球磁场相互作用，将带电粒子引导至地球高层大气，其相互作用形成极光（aurora），可在地磁极附近观测到。\(^\text{[95]}\)日球层中最大的稳定结构是 日球电流片（heliospheric current sheet），它是由太阳自转磁场作用于行星际介质形成的螺旋形结构。\(^\text{[96][97]}\)
\subsection{内太阳系}
内太阳系是指包含类地行星与小行星的区域。[98]该区域内的天体主要由硅酸盐和金属组成，[99]并且它们距离太阳相对较近；整个区域的半径小于木星与土星轨道之间的距离。此区域位于冰霜线（frost line）之内，该冰霜线距离太阳略小于 5 个天文单位（AU）。[52]
\subsubsection{类地行星}
\begin{figure}[ht]
\centering
\includegraphics[width=8cm]{./figures/d502b9668b686c61.png}
\caption{四颗类地行星：水星、金星、地球与火星（The four terrestrial planets Mercury, Venus, Earth and Mars）} \label{fig_Solar_10}
\end{figure}
四颗类地行星（或称内行星）具有致密的岩质组成，卫星稀少或完全没有，也不存在行星环系统。它们主要由耐高温的矿物（如硅酸盐）组成，这些矿物形成了行星的地壳与地幔；而其核心则主要由铁与镍等金属构成。四颗行星中的三颗——金星、地球和火星——拥有足够厚的气体包层，能够形成天气系统；所有类地行星都具有陨石坑和构造地貌，如裂谷与火山。[100]
\begin{itemize}
\item 水星（距离太阳 0.31–0.59 AU）[D 6]是太阳系中最小的行星。其表面呈灰色，拥有由逆冲断层形成的广阔断崖（rupes）系统，以及由撞击事件残迹形成的明亮射纹系统。[101]表面温度变化剧烈——赤道区域夜间温度可达 −170 °C（−270 °F），而白天则高达 420 °C（790 °F）。在过去，水星曾有火山活动，形成了类似月球的光滑玄武质平原。[102]水星可能拥有硅酸盐地壳与巨大的铁核。[103][104]它的大气极其稀薄，仅由太阳风粒子与喷射出的原子组成。[105]水星没有天然卫星。[106]
\item 金星（距离太阳 0.72–0.73 AU）[D 6]拥有高度反射的白色大气，主要由二氧化碳组成。其表面气压约为地球海平面气压的 90 倍。[107]金星表面温度超过 400 °C（752 °F），主要原因是其大气中温室气体含量极高。[108]该行星缺乏可抵御太阳风剥蚀的磁场，这表明其大气可能由火山活动持续补充。[109]其表面显示出广泛的火山活动证据，属于“停滞盖板构造”（stagnant-lid tectonics）。[110]金星没有天然卫星。[106]
\item 地球（距离太阳 0.98–1.02 AU）[D 6]是宇宙中目前已知唯一存在生命与地表液态水的天体。[111]地球大气中含有 78\% 的氮气与 21\% 的氧气，这正是生命活动的结果。[112][113]地球拥有复杂的气候与天气系统，不同气候带之间差异显著。[114]地表主要由绿色植被、荒漠与白色冰盖组成。[115][116][117]地球的板块构造塑造了大陆形态。[102]其行星磁层（magnetosphere）能屏蔽辐射，减少大气流失，从而维持生命宜居性。[118]
\item 月球是地球唯一的天然卫星。[119]其直径约为地球的四分之一。[120]表面覆盖着极细的风化层（regolith），并布满撞击坑。[121][122]月球上的大暗斑——“月海”（maria）——由古代火山活动形成。[123]其大气极其稀薄，仅为部分真空，粒子密度低于每立方厘米 10^7。[124]
\item 火星（距离太阳 1.38–1.67 AU）[D 6]其半径约为地球的一半。[125]由于火星土壤中的氧化铁，其表面呈红色。[126]极地地区覆盖着由水和二氧化碳组成的白色冰盖。[127]火星大气主要由二氧化碳组成，表面气压仅为地球的 0.6\%，但仍足以产生一些天气现象。[128]火星年为 687 个地球日，其表面温度变化剧烈，范围从 −78.5 °C（−109.3 °F）至 5.7 °C（42.3 °F）。表面布满火山与裂谷，并富含多种矿物。[129][130]火星内部结构高度分化，并在约 40 亿年前失去了磁层。[131][132]火星拥有两颗微小卫星：[133]
\item 火卫一（Phobos）是火星的内卫星。其形状不规则，平均半径为 11 km（7 mi）。表面反射率极低，布满撞击坑。[D 7][134]其中最大的撞击坑为 Stickney，半径约 4.5 km（2.8 mi）。[135]
\item 火卫二（Deimos）是火星的外卫星。与火卫一相似，形状不规则，平均半径 6 km（4 mi），表面反光度同样极低。[D 8][D 9]然而，火卫二的表面比火卫一更为光滑，因为风化层部分覆盖了撞击坑。[136]
\end{itemize}
\subsubsection{小行星}
\begin{figure}[ht]
\centering
\includegraphics[width=8cm]{./figures/e8184c45bcc8665a.png}
\caption{} \label{fig_Solar_11}
\end{figure}
除最大的小行星谷神星（Ceres）外，其他小行星均被归类为“太阳系小天体”（small Solar System bodies）。它们主要由碳质物质、耐高温的岩石与金属矿物组成，部分含有冰。[137][138]其尺寸范围从数米至数百公里不等。许多小行星依据其轨道特征被划分为不同的小行星族群（groups）与家族（families）。部分小行星拥有绕其运行的天然卫星，即所谓“次级小行星”——绕更大行星体运行的小行星。[139]
\begin{itemize}
\item 穿越水星轨道的小行星（Mercury-crossing asteroids）指近日点轨道位于水星轨道以内的小行星。截至目前已知至少有 362 颗，它们包括太阳系中距离太阳最近的天体。[140]目前尚未发现任何位于水星轨道与太阳之间的“火神星族”（vulcanoids）。[141][142]截至 2024 年，仅发现一颗完全在金星轨道以内运行的小行星——594913ꞌAylóꞌchaxnim。[143]
\item 穿越金星轨道的小行星（Venus-crossing asteroids）指轨道与金星轨道相交的小行星。截至 2015 年，共记录 2,809 颗。[144]
\item 近地小行星（Near-Earth asteroids）其轨道相对接近地球轨道。[145]部分被列为潜在危险天体（potentially hazardous objects），因为它们未来有可能与地球发生碰撞。[146][147]截至 2024 年，已知的近地小行星数量超过 37,000 颗。[148]一些绕太阳运行的流星体（meteoroids）在撞击地球前足够大，能够被直接观测与追踪。如今普遍认为，过去的撞击事件在塑造地球的地质与生物演化历史中起到了重要作用。[149]
\item 穿越火星轨道的小行星（Mars-crossing asteroids）指近日点大于 1.3 AU、轨道穿越火星轨道的小行星。[150]截至 2024 年，美国 NASA 已确认 26,182 颗火星交叉小行星。[144]
\end{itemize}




